\subsection{primitive 两步盲性、二维加权 \texorpdfstring{$\zeta$}{zeta} 的零温原子化与 resonance window 的 adelic 分裂}\label{subsec:conclusion-primitive-weightedzeta-groundstate-hankel-adelic-splitting}
沿用小节 \ref{subsec:conclusion-atomic-euler-ghost-primitive-core-zerotemp}、小节 \ref{subsec:conclusion-realinput40-uv-atomcore-phase-shell-trace}、小节 \ref{subsec:conclusion-realinput40-analytic-window-logshoulder-ground-quartic-unification}、定理 \ref{thm:conclusion-projective-pressure-integer-algebraic-anchor} 与定理 \ref{thm:conclusion-resonance-window-hankel-null-principalization} 的记号。前述链条已经分别把 primitive/core 的单原子手术、二维加权 \(\zeta(z;u,v)\) 的 atom/core 分裂、输出位势零温地基层、整数阶射影压力锚点以及 resonance window 的 Hankel-null 主理想化固定下来。把这些接口合并之后，还可以继续抽出下面几条此前未被单独封装的终端结论。

\begin{theorem}[primitive 谱对两步 branch 扇区的余一维失明]\label{thm:conclusion-primitive-two-step-codimone-blindness}
设
\[
P(z,u)=\sum_{n\ge 1}p_n(u)z^n,
\qquad
P_{\mathrm{core}}(z,u)=\sum_{n\ge 1}p_n^{\mathrm{core}}(u)z^n
\]
为 full/core primitive 生成函数，并令
\[
\mathcal D:=
\operatorname{span}_{\QQ(u)}
\bigl\{
(p_n(u)-p_n^{\mathrm{core}}(u))_{n\ge 1}
\bigr\}
\subset \QQ(u)^{\NN_{\ge 1}}.
\]
记 \(e_2\) 为仅在 \(n=2\) 坐标非零的标准基向量。则
\[
\mathcal D=\QQ(u)\cdot e_2.
\]
等价地，对任意有限支撑线性泛函
\[
\Lambda=\sum_{n=1}^{N}a_n\,\mathrm{ev}_n
\]
都有
\[
\Lambda(P)-\Lambda(P_{\mathrm{core}})=a_2\,u.
\]
因此，所有满足 \(a_2=0\) 的 primitive 统计量对 full/core 的 branch 差异完全失明。
\leanverified{paper\_conclusion\_primitive\_two\_step\_codimone\_blindness}
\end{theorem}
\begin{proof}
由推论 \ref{cor:conclusion-realinput40-even-defect-polylog-ghost}，
\[
p_n(u)-p_n^{\mathrm{core}}(u)=u\,\mathbf 1_{\{n=2\}}.
\]
故全部差分序列都支撑在唯一坐标 \(n=2\) 上，并由 \(u e_2\) 生成，遂
\[
\mathcal D=\QQ(u)\cdot e_2.
\]
对线性泛函 \(\Lambda\) 逐项配对即得
\[
\Lambda(P)-\Lambda(P_{\mathrm{core}})
=
\sum_{n=1}^{N}a_n\bigl(p_n(u)-p_n^{\mathrm{core}}(u)\bigr)
=
a_2u.
\]
证毕。
\end{proof}

\begin{theorem}[二维联合势的局部一致原子化]\label{thm:conclusion-realinput40-twoparam-local-uniform-atomicization}
设
\[
\Delta(z;u,v)=\det(I-zM_{u,v}),
\qquad
\zeta(z;u,v)=\Delta(z;u,v)^{-1}.
\]
令
\[
z=\sqrt{w/u},
\qquad
u\to+\infty.
\]
则对任意紧集
\[
K\Subset\{(w,v)\in(0,\infty)^2:\ vw\neq 1\}
\]
都有局部一致估计
\[
\sup_{(w,v)\in K}
\left|
\Delta\!\left(\sqrt{w/u};u,v\right)-(1-vw)
\right|
=
O_K(u^{-1/2}),
\]
以及
\[
\sup_{(w,v)\in K}
\left|
\zeta\!\left(\sqrt{w/u};u,v\right)-\frac{1}{1-vw}
\right|
=
O_K(u^{-1/2}).
\]
故在碰撞零温缩放下，二维联合势并非仅点态地，而是局部一致地退化为单一长度 \(2\) 的 Euler 原子。
\end{theorem}
\begin{proof}
定理 \ref{thm:xi-time-part73-twoparam-collision-freezing-single-atom-collapse} 已给出点态展开
\[
P_8\!\left(\sqrt{w/u};u,v\right)=-1+vw+O(u^{-1/2}),
\qquad
\Delta\!\left(\sqrt{w/u};u,v\right)=1-vw+O(u^{-1/2}).
\]
而把显式八次式 \(P_8(z;u,v)\) 代入 \(z=\sqrt{w/u}\) 后，余项是有限个
\[
u^{-a/2}w^bv^c
\]
型单项式之和；故在任意固定紧集 \(K\) 上，各系数一致有界，从而误差可提升为 \(O_K(u^{-1/2})\) 的局部一致估计。

又因 \(K\) 与超曲面 \(vw=1\) 正距分离，存在 \(c_K>0\) 使
\[
|1-vw|\ge c_K
\qquad
((w,v)\in K).
\]
于是对充分大的 \(u\)，\(\Delta(\sqrt{w/u};u,v)\) 在 \(K\) 上亦有统一正下界，取倒数便得到
\[
\zeta\!\left(\sqrt{w/u};u,v\right)
=
\frac{1}{1-vw}+O_K(u^{-1/2}).
\]
证毕。
\end{proof}

\begin{theorem}[输出位势切片的九模律与碰撞冻结角的七模律]\label{thm:conclusion-realinput40-v1-nine-mode-freezecorner-seven-mode}
在真实输入 \(40\) 状态核上，沿输出位势切片 \(v=1\) 有
\[
\deg_z\Delta(z;u,1)\le 9,
\qquad
\dim\ker M_{u,1}\ge 31.
\]
在碰撞冻结角 \((u,v)=(0,1)\) 处更有
\[
\deg_z\Delta(z;0,1)\le 7,
\qquad
\dim\ker M_{0,1}\ge 33.
\]
并且非刚性迹递推的阶数分别至多为 \(7\) 与 \(5\)。因此，在输出位势切片上，\(40\) 维真实输入核的全部非零谱内容至多由 \(9\) 个模态承载；在冻结角上则进一步塌缩到至多 \(7\) 个模态。
\end{theorem}
\begin{proof}
这正是定理 \ref{thm:xi-time-part73-v1-rank-collapse-codim2-tip} 的直接重述。证毕。
\end{proof}

\begin{theorem}[一参数零温极限保留正熵地基，而二维零温极限彻底原子化]\label{thm:conclusion-realinput40-oneparam-ground-vs-twoparam-atomic-collapse}
将一参数输出位势写作
\[
\Delta_{\mathrm{out}}(z;u_{\mathrm{out}})
=(1-u_{\mathrm{out}}z^2)\,F(z;u_{\mathrm{out}}).
\]
令 \(w=u_{\mathrm{out}}z^2\) 并令 \(u_{\mathrm{out}}\to+\infty\)。则 core 部分存在极限
\[
\widehat F(w)
:=
\lim_{u_{\mathrm{out}}\to+\infty}F(\sqrt{w/u_{\mathrm{out}}};u_{\mathrm{out}})
=
\det(I-wB_\infty),
\]
其中
\[
B_\infty=
\begin{pmatrix}
0&1&1&0\\
0&0&1&0\\
0&0&3&1\\
1&0&0&3
\end{pmatrix},
\]
并且
\[
h_{\mathrm{ground}}
=
\frac12\log\rho(B_\infty)>0.
\]
若再记
\[
b=\rho(B_\infty)^{-1},
\]
则主极点常数 \(C_\infty\) 与 Abel 有限部常数 \(\log\mathfrak M_\infty\) 都由该单一代数参数 \(b\) 完整决定。

与此相反，在二维碰撞--输出联合势中，结论 \ref{thm:conclusion-realinput40-twoparam-local-uniform-atomicization} 表明同类零温缩放极限直接坍缩为
\[
1-vw.
\]
因此，一参数输出位势极限保留了一个正熵的四状态地基 SFT，而二维碰撞--输出极限则把该地基完全从极限行列式中消去，只剩单一 Euler 原子。
\end{theorem}
\begin{proof}
由推论 \ref{cor:conclusion-realinput40-nonatomic-freezing-positive-entropy}，
\[
\Delta\!\left(\sqrt{w/u_{\mathrm{out}}},u_{\mathrm{out}}\right)
\longrightarrow
(1-w)\det(I-wB_\infty),
\qquad
h_{\mathrm{ground}}=\frac12\log\rho(B_\infty)>0.
\]
再由定理 \ref{thm:conclusion-realinput40-zero-temp-quartic-constant-unification}，零温常数
\[
b=\rho(B_\infty)^{-1}
\]
同时统一控制 \(h_{\mathrm{ground}}, C_\infty\) 与 \(\log\mathfrak M_\infty\)。另一方面，二维缩放极限由定理 \ref{thm:conclusion-realinput40-twoparam-local-uniform-atomicization} 给出。将两者并置便得结论。证毕。
\end{proof}

\begin{theorem}[resonance window 的 Hankel-null 模完全主理想化]\label{thm:conclusion-resonance-window-hankel-null-principal-module}
对 audited resonance window \(q=9,\dots,17\)，记 \(P_q(\lambda)\) 为审计过的极小特征多项式，\(H^{\mathsf T}\) 为相应 Hankel 主块，并令
\[
\Delta_{\mathrm{prim}}(q):=
\dim\!\bigl(\ker(H^{\mathsf T})/M(P_q)\bigr),
\]
其中 \(M(P_q)\) 表示由 \(P_q\) 生成的倍乘子模。则
\[
\Delta_{\mathrm{prim}}(q)=0
\qquad
(q=9,\dots,17),
\]
从而
\[
\ker(H^{\mathsf T})=M(P_q)
\qquad
(q=9,\dots,17).
\]
换言之，在该审计窗口内不存在任何 primitive Hankel-null 方向；全部 null 证书都只是极小多项式的倍乘冗余。
\end{theorem}
\begin{proof}
这正是定理 \ref{thm:conclusion-resonance-window-hankel-null-principalization} 的 primitive quotient 表述。证毕。
\end{proof}

\begin{corollary}[Hankel-null 证书对极小递推的 gcd 恢复律]\label{cor:conclusion-resonance-window-hankel-null-gcd-recovery}
对 \(q=9,\dots,17\)，任取一组生成 \(\ker(H^{\mathsf T})\) 的整数 Hankel-null 基，并把相应零模写成整系数湮灭多项式族 \(\{B_\alpha(\lambda)\}\)。则
\[
\gcd_{\ZZ[\lambda]}\{B_\alpha(\lambda)\}=\pm P_q.
\]
等价地，audited resonance window 中的任意整数零模都只是 \(P_q\) 的截断倍乘冗余，而由任意一组 Hankel-null 基都可唯一恢复 \(\pm P_q\)。
\end{corollary}
\begin{proof}
由定理 \ref{thm:conclusion-resonance-window-hankel-null-principal-module}，\(\ker(H^{\mathsf T})=M(P_q)\)。因此每个整数零模都属于由 \(P_q\) 生成的倍乘子模，故对应湮灭多项式都含有 \(P_q\) 因子；反过来，该模本身又由 \(P_q\) 生成，因而任何生成集的整系数最大公因子只能是 \(\pm P_q\)。证毕。
\end{proof}

\begin{theorem}[离散碰撞谱的双实现统一]\label{thm:conclusion-collision-spectrum-dual-realization-unification}
设 \(\mathcal L_q\) 为定义在单纯形 \(\Delta\) 上的射影转移算子族。则对每个整数 \(q\ge 2\)，其在齐次次数 \(q\) 子空间 \(\mathcal H_q\) 上的限制满足
\[
\mathcal L_q|_{\mathcal H_q}=\widetilde A_q,
\qquad
\rho(\mathcal L_q|_{\mathcal H_q})=\rho(\widetilde A_q)=r_q.
\]
另一方面，在 arity--charge 一参数切片
\[
u=q^{-1},
\qquad
v=q
\]
上，二维加权行列式退化为
\[
\det(I-zM(q))=(z+1)(1-qz^2)Q_7(z,q).
\]
因此，同一整数参数 \(q\) 同时精确控制两套表面上不同的对象：一套是射影动力学中的对称幂 moment-kernel，另一套是 weighted primitive-\(\zeta\) 理论中的低阶行列式切片。故离散碰撞谱并非孤立的整数序列，而是同一 \(q\)-参数热力学家族在两个实现范畴中的同步截面。
\end{theorem}
\begin{proof}
第一部分即定理 \ref{thm:conclusion-projective-pressure-integer-algebraic-anchor}。第二部分即定理 \ref{thm:conclusion-realinput40-arity-charge-z2-trace-splitting} 的三因子分解。将二者并置，可见整数参数 \(q\) 在两条独立构造链上共享同一离散谱锚点 \(r_q\) 与同一低阶显式因子 \(1-qz^2\)。证毕。
\end{proof}

\begin{theorem}[resonance window 中的 Archimedean--模 \(2\) 复杂度分裂]\label{thm:conclusion-resonance-window-archimedean-mod2-complexity-splitting}
对 \(q=9,\dots,17\)，collision moment 序列 \(S_q(m)\) 同时满足如下二元刚性：
\begin{enumerate}
\item 在模 \(2\) 上，经有限平移 \(m\mapsto m+a_q\) 后，
\[
\Delta^{b_q}S_q(m+a_q)\equiv 0\pmod 2,
\qquad
\Delta:=E+1,
\]
从而 \(S_q(m+a_q)\bmod 2\) 在后段塌缩为次数 \(<b_q\) 的离散二项式多项式，并且其最终周期只落在三种 dyadic 级别
\[
T_q\mid 4,\qquad T_q\mid 8,\qquad T_q\mid 16.
\]
\item 在 Archimedean 侧，同一 \(q\)-序列仍由 Perron 根 \(r_q=\rho(\widetilde A_q)\) 控制，并保持指数增长尺度
\[
S_q(m)=\Theta(r_q^m),
\qquad
r_q>1.
\]
\end{enumerate}
因此，在同一 resonance window 内，有限特征阴影的复杂度坍缩为极低周期，而实数域侧却维持强指数增长谱尺度；二者并不同步退化。更强地，不同的 \(q\) 甚至可以拥有完全相同的模 \(2\) 影子，却具有不同的 Archimedean 越线强度。
\end{theorem}
\begin{proof}
模 \(2\) 部分由定理 \ref{thm:conclusion-resonance-window-mod2-binomial-collapse} 与定理 \ref{thm:conclusion-resonance-window-dyadic-chambers} 给出。Archimedean 部分中，
\[
r_q=\rho(\widetilde A_q)
\]
由定理 \ref{thm:conclusion-projective-pressure-integer-algebraic-anchor} 给出，它正是整数阶 pressure 锚点的指数尺度；于是同一 \(q\) 的实数域增长仍按 \(r_q\) 指数级控制。最后，模 \(2\) 影子不足以决定 Archimedean 不稳定强度，则由推论 \ref{cor:conclusion-mod2-shadow-insufficient-for-archimedean-instability} 直接给出。证毕。
\end{proof}

\begin{conjecture}[黄金协议下剩余障碍的几何--时间封口]\label{conj:conclusion-golden-protocol-residual-geometry-time-closure}
沿用定理 \ref{thm:conclusion-resonance-window-hankel-null-principal-module} 与推论 \ref{cor:conclusion-resonance-window-hankel-null-gcd-recovery} 的 resonance-window 口径。若再把协议端点固定到定理 \ref{thm:conclusion-protocol-acceptable-golden-delay-rigidity} 的黄金刚性端点，则任何仍无法被 strictify、可逆化或单值审计化的剩余障碍，都应当只来自两类机制：其一是定理 \ref{thm:conclusion-median-prime-support-minimal-rigidity} 所刻画的 \(\Theta\)-类/prime-support 几何不可压缩性；其二是推论 \ref{cor:conclusion-delay3-not-absorbable-by-static-ledger} 所刻画的 \(\delta_{\min}=3\) 同步时间残差。换言之，一旦标量递推谱已由 \(P_q\) 及其 Hankel-null 主理想化完全封口，剩余不可约内容便不再驻留于一维递推数据，而只能驻留于几何可见轴数与在线时间量子之中。
\end{conjecture}

\begin{conjecture}[fold collision 层级的普适 adelic 不对称律]\label{conj:conclusion-fold-collision-universal-adelic-asymmetry}
设 resonance window 中观察到的两步饱和机制在全部高阶 \(q\) 上持续，即
\[
r_q^{1/q}\to \sqrt{\varphi},
\qquad
\eta_q:=\frac{|\mu_q|}{r_{q-2}}\to 1.
\]
则应存在如下普适分裂：
\begin{enumerate}
\item \textbf{Archimedean 侧：}
\[
q\,g_q^{(c)}\to 2,
\qquad
\delta_{\mathrm{ait}}(q)\to \varphi^{-1},
\]
即复缺口仅以 \(2/q\) 级缓慢闭合，而 Aitken 主误差趋于黄金倒数。
\item \textbf{非阿基米德侧：}
对每个固定素数 \(p\)，\(p\)-进 shadow 在有限平移后应进入有界状态空间；特别地，\(p=2\) 时这已在 audited resonance window 中表现为周期至多 \(16\) 的离散多项式影子。
\end{enumerate}
若该猜想成立，则 fold collision 层级呈现一种真正的 adelic asymmetry：Archimedean 模态保留普适但非平凡的渐近缺口律，而每个固定 \(p\)-进通道都只留下低复杂度、有限状态的阴影。
\end{conjecture}

综上，primitive--Witt 层的 full/core 差异首先坍缩为唯一的长度 \(2\) 坐标缺口；二维 weighted-\(\zeta\) 在碰撞零温缩放下进一步局部一致地原子化为 \(1-vw\)；但一参数输出位势的零温端却仍保留一个正熵四状态地基层。与此同时，resonance window 内的 Hankel-null 模完全主理想化，并且任意整数 Hankel-null 基都只能通过 gcd 恢复同一个 \(\pm P_q\)；这说明标量递推侧的原始约束在审计窗口内已经完全封口。整数阶射影压力与 arity--charge 行列式切片又共享同一离散谱锚点 \(q\)，而模 \(2\) 阴影的低复杂度坍缩与 Archimedean 谱尺度的持续增长被迫并存，从而把整个 collision hierarchy 推向一种真正的 adelic 分裂。若进一步结合猜想 \ref{conj:conclusion-golden-protocol-residual-geometry-time-closure}，则在黄金协议端点上，标量递推之后的剩余不可约障碍应只剩几何 prime-support 轴数与 \(\delta_{\min}=3\) 时间量子两类来源。

\endinput
